
\chapter{Detalles de la propuesta}
\label{cap:propuesta}

Este capítulo desarrolla la propuesta del trabajo: un prototipo determinista que
detecta y combina condicionales anidados de forma segura. Se presenta primero
una visión general del enfoque y su encaje en la metodología \emph{Design
Science Research}; a continuación, el diseño del enfoque, con el patrón objetivo
y las precondiciones de aplicabilidad que delimitan cuándo la transformación es
segura; después, la arquitectura del prototipo en paquetes y flujos de datos; y,
por último, los detalles de implementación de cada componente.

\section{Visión general}
\label{sec:vision-general}
La propuesta se plantea como un proyecto de \emph{Design Science Research}
(DSR)~\cite{10.2307/25148625, 10.2753/MIS0742-1222240302}, porque su
estructura encaja de forma natural con el
objetivo: construir un artefacto y evaluarlo. El artefacto es un prototipo
determinista de detección y refactorización; la evaluación consiste en
medir cómo responde a cada una de las preguntas de investigación. Esta
doble naturaleza (construir y comprobar) recorre todo el trabajo. Conviene
subrayar, además, que el prototipo no es solo el objeto de estudio: actúa
como \emph{baseline verificable}, es decir, como la referencia fiable
contra la que se contrastan después las refactorizaciones de los grandes
modelos de lenguaje en RQ3.

\paragraph{Mapa RQ $\leftrightarrow$ secciones.}
La \autoref{tab:mapa-rq} relaciona cada pregunta de investigación con las
secciones que la abordan.

\begin{table}[ht]
\centering
\caption{Trazabilidad entre preguntas de investigación y secciones.}
\label{tab:mapa-rq}
\begin{tabular}{@{}lll@{}}
\toprule
\textbf{RQ} & \textbf{Aborda} & \textbf{Secciones} \\
\midrule
RQ1 & Casos combinables y resultado & \ref{sec:diseno}, \ref{sec:precondiciones}, \ref{sec:resultados-rq1} \\
RQ2 & Impacto en complejidad cognitiva & \ref{sec:proxy}, \ref{sec:exp-rq2}, \ref{sec:resultados-rq2} \\
RQ3 & LLMs: corrección y automatización & \ref{sec:llm-subsistema}, \ref{sec:exp-rq3}, \ref{sec:resultados-rq3} \\
\bottomrule
\end{tabular}
\end{table}


